\section{Linear autonomous delay systems and the argument principle}\label{sec:theory}

\subsection{System classes and the characteristic function}\label{sec:classes}

Consider the linear autonomous delay system in first-order descriptor form
\begin{equation}\label{eq:ddae}
  \mat E\,\dot{\vect x}(t) = \mat A_0\,\vect x(t)
  + \sum_{j=1}^{m}\Bigl(\mat A_j\,\vect x(t-\tau_j) + \mat B_j\,\dot{\vect x}(t-\tau_j)\Bigr)
  + \int_{0}^{\vartheta_{\mathrm{max}}} \mat W(\vartheta)\,\vect x(t-\vartheta)\,\dd\vartheta ,
\end{equation}
with state $\vect x(t)\in\mathbb R^{N}$, constant matrices
$\mat E,\mat A_j,\mat B_j\in\mathbb R^{N\times N}$, discrete delays
$\tau_j>0$ and a distributed-delay kernel $\mat W$. The exponential ansatz
$\vect x(t)=\vect\xi\,\ee^{\lam t}$ leads to the characteristic matrix and
the characteristic function
\begin{equation}\label{eq:charfun}
  \mat M(\lam) = \lam\mat E - \mat A_0
   - \sum_{j=1}^{m}\bigl(\mat A_j + \lam\mat B_j\bigr)\ee^{-\lam\tau_j}
   - \int_{0}^{\vartheta_{\mathrm{max}}}\mat W(\vartheta)\,\ee^{-\lam\vartheta}\dd\vartheta,
  \qquad
  \Dfun(\lam) = \det\mat M(\lam).
\end{equation}
Equation~\eqref{eq:ddae} covers all the system classes treated in this paper:
\emph{retarded} DDEs ($\mat E=\mat I$, $\mat B_j=\mat 0$), \emph{neutral}
DDEs ($\mat B_j\neq\mat 0$), delay \emph{differential--algebraic} equations
(DDAEs, singular $\mat E$, e.g., due to kinematic constraints or algebraic
force balances), and \emph{distributed} delays. Beyond~\eqref{eq:ddae}, the
method below requires only an analytic function $\Dfun(\lam)$ that can be
evaluated at complex arguments; hence transcendental characteristic functions
of continua (e.g., $\cosh$-type expressions of elastic
beams~\cite{zhang2016exact}) and fractional-order quasi-polynomials
containing terms $\lam^{\alpha}$, $\alpha\notin\mathbb N$,
are admissible as well.

The equilibrium of~\eqref{eq:ddae} is asymptotically stable if all roots of
$\Dfun(\lam)=0$ satisfy $\Real\lam<0$. More generally, a guaranteed decay
rate (spectral gap) $|\sigma|$ -- often called $\gamma$-stability -- requires
all roots to lie left of the shifted line $\Real\lam=\sigma<0$. We denote by
$\Zint(\vect p)$ the number of roots with $\Real\lam>\sigma$ for the design
parameters $\vect p$; the stability chart is the map
$\vect p\mapsto\Zint(\vect p)$, and the stable (or $\gamma$-stable) domain
is the region $\Zint=0$.

For large $|\lam|$ the characteristic function grows like
$\Dfun(\lam)\sim c\,\lam^{n}$, where the leading order $n$ equals $N$ for
retarded systems in state-space form but is smaller for DDAEs (the algebraic
rows remove dynamic content) and is only an \emph{effective, possibly
non-integer} order for neutral ($\Dfun\sim\lam^{n}(1+\sum_j b_j
\ee^{-\lam\tau_j})$) and fractional systems. The counting formula below needs
$n$, and Section~\ref{sec:phaseode} shows that it can be estimated
numerically without any structural knowledge.

\paragraph{Fractional orders and the branch cut.} Fractional powers deserve an
explicit word here, because they are the one class in which the choice of
integration line is a question of validity rather than of cost. The principal
branch of $\lam^{\mu}$ is analytic on the complex plane cut along the negative
real axis, and $\Dfun$ inherits that cut: for the fractional controller of
Appendix~\ref{app:gallery} the two sides of the cut differ by
$|\Dfun(-0.25+\ii0^{+})-\Dfun(-0.25-\ii0^{-})|=\FracJump$, so $\Dfun$ is
\emph{discontinuous} there, not merely non-smooth.

The cut lies on the \emph{negative} reals, hence entirely outside the closed
right half-plane, and the imaginary axis itself is untouched: the same
measurement across the axis returns a jump of $\FracAxisJump$ -- identically
zero -- at every frequency tested. The choice $\sigma=0$ is therefore
admissible, and this is not a technicality about one example: it is what makes
the whole class available to the method. Two provisos come with it. First, the
branch \emph{point} at $\lam=0$ lies \emph{on} the contour, where
$\Dfun'\sim\mu\lam^{\mu-1}$ diverges for $\mu<1$; the resulting integrand
singularity $\sim\omega^{\mu-1}$ is integrable for $\mu>0$, and the standard
indentation of the contour around the origin contributes nothing in the limit
\emph{provided} $\Dfun(0)\neq0$. Second, when $\Dfun(0)=0$ the origin is itself
a characteristic root sitting on the contour and no count is defined -- for the
fractional PI controller $\Dfun(0)=k_i$, so the line $k_i=0$ crossing that
chart is exactly such a degeneracy, and the package reports it through the
invalid marker of Section~\ref{sec:integerness} rather than by inventing a
number. That the construction then works is checkable rather than assumed: over
$\FracResidN$ points of the two fractional charts the integer residual has
median $\FracResidMed$ and never exceeds $\FracResidMax$.

A \emph{shifted} line ($\sigma<0$) is a different matter entirely: it encloses
the segment $[\sigma,0]$ of the cut, $\Dfun$ is not analytic inside the
contour, and the argument principle does not apply. Appendix~\ref{app:gallery}
demonstrates what that costs.

\subsection{Counting unstable roots on a half-line}\label{sec:counting}

Let the integration line $\lam=\sigma+\ii\omega$ be free of characteristic
roots, and let $\Dfun$ be analytic with finitely many roots to the right of the
line. Applying the argument principle to the D-shaped contour composed of the
line and the arc at infinity, and exploiting the conjugate symmetry
$\Dfun(\overline\lam)=\overline{\Dfun(\lam)}$ of real-coefficient systems,
the number of roots to the right of the line is obtained from frequencies
$\omega\ge 0$ only~\cite{stepan1989retarded,hassard1997counting}:
\begin{equation}\label{eq:count}
  \Zint \;=\; \frac{n}{2} \;-\; \frac{1}{\pi}\,\Delta\theta,
  \qquad
  \Delta\theta = \int_{0}^{\infty}\theta'(\omega)\,\dd\omega,
  \qquad
  \theta(\omega) = \arg \Dfun(\sigma+\ii\omega),
\end{equation}
where $\theta$ is continuous in $\omega$ and $n/2$ accounts for the phase
gathered on the arc at infinity, where $\Dfun\sim c\lam^{n}$. This is the
half-line (Mikhailov-type) form of the Nyquist criterion. Writing
$\Dfun'(\omega)\equiv\partial_\omega\Dfun(\sigma+\ii\omega)$ for the
derivative along the line, the phase slope is
\begin{equation}\label{eq:integrand}
  \theta'(\omega)
  = \Imag\frac{\Dfun'(\omega)}{\Dfun(\sigma+\ii\omega)}
  = \frac{\Real\Dfun\cdot\partial_\omega\Imag\Dfun
        - \Imag\Dfun\cdot\partial_\omega\Real\Dfun}{|\Dfun|^{2}} .
\end{equation}
For $\sigma=0$, $\Zint$ counts the unstable roots; for $\sigma<0$ it counts
the roots violating the prescribed spectral gap, so the same formula draws
$\gamma$-stability charts.

Three standing assumptions deserve explicit statement. \emph{(i)~Finitely many
roots right of the line, none on it.} For retarded systems and regular DDAEs
of low index this holds for every $\sigma$ greater than some abscissa; for
DDAEs of higher index the spectrum can acquire advanced-type root chains that
run into the right half-plane, and the count is then infinite for any
line~\cite{michiels2011spectrum,ha2016analysis} -- such systems are outside
the scope of~\eqref{eq:count}, although the diagnostic of
Section~\ref{sec:integerness} flags them, because the truncated integral then
fails to settle to an integer. \emph{(ii)~Neutral tail.} For neutral systems
$\Dfun\sim c\lam^{n}\bigl(1+\sum_j b_j\ee^{-\lam\tau_j}\bigr)$ and the phase
of the parenthesis oscillates without decay as $\omega\to\infty$: the
integral in~\eqref{eq:count} does not converge but oscillates within a
bounded envelope. If the neutral amplitude satisfies $\beta=\sum_j|b_j|<1$ --
the strong-stability condition of the difference
operator~\cite{hale2002strong,michiels2005eigenvalue} -- that envelope is
bounded by $\arcsin(\beta)/\pi<1/2$ per delay term, so truncating at a
generic $\wmax$ and rounding still recovers the exact count. As $\beta\to1$
(essential instability approaching) the bound saturates and the integer
residual of Section~\ref{sec:integerness} grows accordingly -- the diagnostic
degrades exactly where the answer itself becomes fragile, which is the honest
behaviour (see the acceleration-feedback case study A.6 in
Appendix~\ref{app:gallery}). \emph{(iii)~Analyticity right of the line.} Fractional
powers $\lam^{\alpha}$ are taken on the principal branch, whose cut lies on
the negative real axis; for $\sigma\ge0$ the region right of the line is then
free of branch points, and~\eqref{eq:count} applies verbatim. A shifted line
with $\sigma<0$, however, crosses the cut, and the count refers to the roots
of the principal-branch extension; case study A.11 in
Appendix~\ref{app:gallery} discusses the practical consequences.

\subsection{Why the naive evaluation is impractical}\label{sec:naive}

Formula~\eqref{eq:count} looks like a textbook quadrature problem, but three
properties make its direct use for stability charts inefficient.
First, $\Zint(\vect p)$ is integer-valued: it is constant inside every cell
of the chart and jumps at the boundaries, so the map cannot be interpolated,
and any grid-based chart must be fine enough to resolve every cell of
interest. Second, the numerical evaluation is hardest exactly where it
matters most. If a root $\lam_r=\sigma_r+\ii\omega_r$ lies close to the
integration line, the integrand~\eqref{eq:integrand} develops a
near-Lorentzian peak at $\omega\approx\omega_r$,
\begin{equation}\label{eq:peak}
  \theta'(\omega) \approx \frac{-(\sigma_r-\sigma)}
  {(\sigma_r-\sigma)^{2} + (\omega-\omega_r)^{2}} + \text{smooth part},
\end{equation}
whose half-width $|\sigma_r-\sigma|$ tends to zero as the parameter point
approaches the stability boundary, while its integral contribution stays
$\pm\pi$; \emph{any} method that samples the integrand pointwise and happens
to place all its samples outside such a peak miscounts $\Zint$ by exactly
one, and -- since the result stays an integer -- does so without any
warning. Third, the integration range is infinite: the phase
slope decays only algebraically, with delay-induced oscillatory ripples, so a
fixed truncation $\wmax$ must balance truncation error against the cost of
resolving a ripple over a huge interval.

Section~\ref{sec:method} addresses the first and third difficulties with the
step-size control of an adaptive integrator, and turns the second into an
asset by extracting the location of the offending root from the very peak
that makes the integration hard. The peak-skipping hazard, by contrast,
cannot be legislated away by any choice of integrator; what can be done --
and what Section~\ref{sec:integerness} does -- is to detect it, using the
fact that the root estimate and the phase count come from independent
mechanisms and must agree.